\section{Conclusion}
\label{sec:conclusion}

Agents read far more than they use, and what they use can be learned from their own traces: on four of the five
agents studied, a projection learned from other runs removed \PhaseAPrudentSavingi{} and \PhaseAPrudentSavingiii{} of
the tool-output tokens of a new run on two agents with real APIs, and \PhaseAPrudentSavingiv{} and
\PhaseAPrudentSavingv{} on two simulated ones (re-serialisation alone, which removes no field, accounting for up to
\ReserialAi{} and \ReserialAiv{}). Where it failed, it failed for a reason: the agent had transformed a value, used a field only
to reason, or, deprived of an identifier, invented one. Two safeguards, keeping identifying fields and recognising
transformed values, removed the losses that motivated them while keeping a saving of \StudyTwoVtwoTokensi{} and
\StudyTwoVtwoTokensiii{} on those real-API agents and \StudyTwoVtwoTokensiv{} and \StudyTwoVtwoTokensv{} on the
simulated ones; a fourth mechanism, translation, remains open.
The method is therefore best used as it is implemented: as a proposal learned from complete traces, reviewed by a
person, applied with identifiers protected and with a way back for the agent when a cut was wrong.
